

\section{Reproducibility Statement}

We have included the heterogeneous data construction process in Section~\ref{sec:task_definition} and more details in the Appendix. We describe the RL training settings and experiment platforms in Appendix~\ref{apd:training_details} and Section~\ref{sec:experimental_setup}. The prompts involving the usage of LLM (primarily GPT-4.1) are provided in Appendix~\ref{apd:prompts}. The above materials should be sufficient to reproduce our work.



\section{Data Process Details} \label{apd:data}

\subsection{Detailed Introduction of Datasets}
\noindent\textbf{Translation (En-Fr, Fr-En)}: This task requires the LLM to translate between English and French (in our case, English to French and French to English). We use the dataset \texttt{aircrypto/English-French-Translations-Train-Large}~\citep{hf_enfr} from HuggingFace, which provides high-quality, paired sentence-level samples.

\noindent\textbf{Instruction Following}: Given specific requirements in the provided instructions, the LLM should respond accordingly. We use \texttt{tulu3-sft-reused-on-policy-8b}, part of the Tulu-3~\citep{tulu3} preference dataset, which contains generation pairs between different LLMs during the training of Llama‑3.1‑Tulu‑3‑8B.

\noindent\textbf{Multi-turn}: LLMs respond to instructions with previous interaction histories. We pick \texttt{allenai-WildChat-1M-multiturn}~\citep{allenwilde}, a collection of 1M ChatGPT interaction logs from the wild. We select the English subset aimed at RLHF queries. 

\noindent\textbf{Summarization}: This task requires the LLM to summarize long documents into short abstracts. We pick \texttt{ccdv/govreport-summarization, ccdv/pubmed-summarization, ccdv/arxiv-summarization}~\citep{cohan-etal-2018-discourse-summarizationdataset}, which include different types of documents from arXiv articles to government reports.

\noindent\textbf{Math}: We pick OpenAI \textsc{gsm8k}~\citep{gsm8k}, a classic dataset of grade‑school math problems designed to evaluate multi‑step reasoning. We choose not to use more complex math-reasoning datasets because our focus in this work is primarily on LLM text-generation tasks. Advanced math reasoning often requires specialized methodologies, such as tree‑search reasoning, which makes it unsuitable for single-pass direct generation.

\noindent\textbf{Conditional Generation}: The LLM should generate coherent text according to given constraints. In our setting, we task the LLM with filling in missing paragraphs in an essay or producing a complete essay based on an abstract outline. We use \texttt{qwedsacf/ivypanda-essays}~\citep{ivypanda}, a HuggingFace dataset repository containing long-form essays covering multiple disciplines sourced from the \textit{IvyPanda platform}\footnote{\texttt{https://ivypanda.com/}}.


\subsection{Data Filtering Prompt} ~\label{apd:prompt_for_data_filter}

\begin{tcolorbox}[promptbox, title={DATA\_FILTER\_PROMPT}]
\textbf{Input}: \texttt{task\_samples}, \texttt{quality\_standards}
\tcblower
You are given a set of task samples, each consisting of:
\begin{enumerate}
    \item \textbf{User Query} – the task or request made to the model.
    \item \textbf{Model Response} – the output given by the model.
\end{enumerate}

The samples may come from various task types, including: Translation, Summarization, Math problem solving, RLHF alignment, Conditional text generation, and Multi-turn dialogue.

\textbf{Your Goal}: Identify and select only the samples that did not meet quality standards based on:

\medskip
\textbf{A. Query Quality Issues}:
\begin{itemize}
    \item Ill-formed or incomplete queries; Ambiguous or misleading instructions.
    \item Irrelevant or off-topic requests; Grammatically broken or nonsensical input.
\end{itemize}

\textbf{B. Response Quality Issues}:
\begin{itemize}
    \item Incorrect or factually wrong answers; Incomplete responses.
    \item Poor language quality or incoherent writing.
    \item Hallucinations or made-up facts; Misinterpretation of the query.
\end{itemize}

\textbf{Instructions}:
\begin{enumerate}
    \item For each sample, examine both the query and response.
    \item Mark the sample as \textbf{"Fail"} if either the query quality or the response quality is below standard.
    \item Briefly explain why the sample fails, citing issues in query, response, or both.
    \item Output only the failing samples, in the format:
    \begin{tcolorbox}[colback=white, colframe=gray!20, size=minimal, left=2mm, top=1mm, bottom=1mm]
        \texttt{[Sample ID]} \\
        \texttt{Query: ...} \\
        \texttt{Response: ...} \\
        \texttt{Fail Reason: ...}
    \end{tcolorbox}
\end{enumerate}

\textit{Be strict in applying the criteria — even if only one side (query or response) is substandard, the sample should be considered as failing.}
\end{tcolorbox}